% Options for packages loaded elsewhere
\PassOptionsToPackage{unicode}{hyperref}
\PassOptionsToPackage{hyphens}{url}
\documentclass[
]{article}
\usepackage{xcolor}
\usepackage{amsmath,amssymb}
\setcounter{secnumdepth}{-\maxdimen} % remove section numbering
\usepackage{iftex}
\ifPDFTeX
  \usepackage[T1]{fontenc}
  \usepackage[utf8]{inputenc}
  \usepackage{textcomp} % provide euro and other symbols
\else % if luatex or xetex
  \usepackage{unicode-math} % this also loads fontspec
  \defaultfontfeatures{Scale=MatchLowercase}
  \defaultfontfeatures[\rmfamily]{Ligatures=TeX,Scale=1}
\fi
\usepackage{lmodern}
\ifPDFTeX\else
  % xetex/luatex font selection
\fi
% Use upquote if available, for straight quotes in verbatim environments
\IfFileExists{upquote.sty}{\usepackage{upquote}}{}
\IfFileExists{microtype.sty}{% use microtype if available
  \usepackage[]{microtype}
  \UseMicrotypeSet[protrusion]{basicmath} % disable protrusion for tt fonts
}{}
\makeatletter
\@ifundefined{KOMAClassName}{% if non-KOMA class
  \IfFileExists{parskip.sty}{%
    \usepackage{parskip}
  }{% else
    \setlength{\parindent}{0pt}
    \setlength{\parskip}{6pt plus 2pt minus 1pt}}
}{% if KOMA class
  \KOMAoptions{parskip=half}}
\makeatother
\usepackage{color}
\usepackage{fancyvrb}
\newcommand{\VerbBar}{|}
\newcommand{\VERB}{\Verb[commandchars=\\\{\}]}
\DefineVerbatimEnvironment{Highlighting}{Verbatim}{commandchars=\\\{\}}
% Add ',fontsize=\small' for more characters per line
\newenvironment{Shaded}{}{}
\newcommand{\AlertTok}[1]{\textcolor[rgb]{1.00,0.00,0.00}{\textbf{#1}}}
\newcommand{\AnnotationTok}[1]{\textcolor[rgb]{0.38,0.63,0.69}{\textbf{\textit{#1}}}}
\newcommand{\AttributeTok}[1]{\textcolor[rgb]{0.49,0.56,0.16}{#1}}
\newcommand{\BaseNTok}[1]{\textcolor[rgb]{0.25,0.63,0.44}{#1}}
\newcommand{\BuiltInTok}[1]{\textcolor[rgb]{0.00,0.50,0.00}{#1}}
\newcommand{\CharTok}[1]{\textcolor[rgb]{0.25,0.44,0.63}{#1}}
\newcommand{\CommentTok}[1]{\textcolor[rgb]{0.38,0.63,0.69}{\textit{#1}}}
\newcommand{\CommentVarTok}[1]{\textcolor[rgb]{0.38,0.63,0.69}{\textbf{\textit{#1}}}}
\newcommand{\ConstantTok}[1]{\textcolor[rgb]{0.53,0.00,0.00}{#1}}
\newcommand{\ControlFlowTok}[1]{\textcolor[rgb]{0.00,0.44,0.13}{\textbf{#1}}}
\newcommand{\DataTypeTok}[1]{\textcolor[rgb]{0.56,0.13,0.00}{#1}}
\newcommand{\DecValTok}[1]{\textcolor[rgb]{0.25,0.63,0.44}{#1}}
\newcommand{\DocumentationTok}[1]{\textcolor[rgb]{0.73,0.13,0.13}{\textit{#1}}}
\newcommand{\ErrorTok}[1]{\textcolor[rgb]{1.00,0.00,0.00}{\textbf{#1}}}
\newcommand{\ExtensionTok}[1]{#1}
\newcommand{\FloatTok}[1]{\textcolor[rgb]{0.25,0.63,0.44}{#1}}
\newcommand{\FunctionTok}[1]{\textcolor[rgb]{0.02,0.16,0.49}{#1}}
\newcommand{\ImportTok}[1]{\textcolor[rgb]{0.00,0.50,0.00}{\textbf{#1}}}
\newcommand{\InformationTok}[1]{\textcolor[rgb]{0.38,0.63,0.69}{\textbf{\textit{#1}}}}
\newcommand{\KeywordTok}[1]{\textcolor[rgb]{0.00,0.44,0.13}{\textbf{#1}}}
\newcommand{\NormalTok}[1]{#1}
\newcommand{\OperatorTok}[1]{\textcolor[rgb]{0.40,0.40,0.40}{#1}}
\newcommand{\OtherTok}[1]{\textcolor[rgb]{0.00,0.44,0.13}{#1}}
\newcommand{\PreprocessorTok}[1]{\textcolor[rgb]{0.74,0.48,0.00}{#1}}
\newcommand{\RegionMarkerTok}[1]{#1}
\newcommand{\SpecialCharTok}[1]{\textcolor[rgb]{0.25,0.44,0.63}{#1}}
\newcommand{\SpecialStringTok}[1]{\textcolor[rgb]{0.73,0.40,0.53}{#1}}
\newcommand{\StringTok}[1]{\textcolor[rgb]{0.25,0.44,0.63}{#1}}
\newcommand{\VariableTok}[1]{\textcolor[rgb]{0.10,0.09,0.49}{#1}}
\newcommand{\VerbatimStringTok}[1]{\textcolor[rgb]{0.25,0.44,0.63}{#1}}
\newcommand{\WarningTok}[1]{\textcolor[rgb]{0.38,0.63,0.69}{\textbf{\textit{#1}}}}
\setlength{\emergencystretch}{3em} % prevent overfull lines
\providecommand{\tightlist}{%
  \setlength{\itemsep}{0pt}\setlength{\parskip}{0pt}}
\usepackage{bookmark}
\IfFileExists{xurl.sty}{\usepackage{xurl}}{} % add URL line breaks if available
\urlstyle{same}
\hypersetup{
  hidelinks,
  pdfcreator={LaTeX via pandoc}}

\author{}
\date{}

\begin{document}

\section{QuasarSTM (Block-STM 3.0) --- Red Adversarial
Review}\label{quasarstm-block-stm-30--red-adversarial-review}

\textbf{Author}: Red (adversarial) \textbf{Date}: 2026-04-26
\textbf{Scope}: LP-010 (QuasarSTM), LP-132 (QuasarGPU), Metal kernel,
layout, host driver, tests \textbf{Target launch}: 2025-12-25
\textbf{Threat model}: An adversary submits transactions designed to
cause two honest validators on the same input to produce divergent
\texttt{state\_root}, \texttt{receipts\_root}, \texttt{execution\_root},
\texttt{block\_hash}, or \texttt{tx\_count}. Any divergence is a chain
split.

\begin{center}\rule{0.5\linewidth}{0.5pt}\end{center}

\subsection{Executive Summary}\label{executive-summary}

\textbf{16 findings: 4 Critical, 5 High, 4 Medium, 3 Low/Informational.}

The 3.0 spec is well thought out --- lanes, three-tier validation,
commit horizons, deterministic policy ordering, and the discipline
against "retry-until-success" loops are real safety primitives.
\textbf{None of that is implemented.} The shipping substrate is closer
to 2.0 with v0.40 DAG bolted on. The gap between the LP and the kernel
is the headline architectural risk: LP-010 documents anti-livelock hard
limits (\texttt{MAX\_FAST\_REPAIRS=3}, \texttt{MAX\_TOTAL\_REPAIRS=8},
\texttt{HOT\_LANE\_CONFLICT\_THRESHOLD=16}) as "mandatory on GPU"; the
kernel enforces none of them. Repair amplification is theoretically
O(N²) and the only thing capping it in practice is the
host\textquotesingle s \texttt{max\_epochs} argument to
\texttt{run\_until\_done}.

The most concerning concrete issue is \textbf{Q3.0-STM-001}: on the slow
path, \texttt{drain\_validate} performs \texttt{mvcc\_apply\_writes}
\emph{before} \texttt{ring\_try\_push(commit\_hdr,\ ...)}. If the commit
ring is full (the spec calls this out as expected backpressure), the
version is bumped, the tx is requeued onto Validate, and on the next
tick it re-validates, finds its own bumped version, and is shoveled to
Repair --- incrementing \texttt{conflict\_count},
\texttt{repair\_count}, and corrupting the \texttt{execution\_root}
chain order across runs because \emph{which} tick the commit ring fills
is timing-dependent. Two validators on the same input will produce
different roots if the host on validator A drains the commit ring with
different latency than validator B.

The next-most-concerning is \textbf{Q3.0-STM-002}: \texttt{mvcc\_locate}
writes \texttt{key\_lo}/\texttt{key\_hi} as plain non-atomic stores. Two
GPU threads racing to claim the same empty slot can both succeed (the
second tx\textquotesingle s store overwrites the first), and a
concurrent reader can observe \texttt{key\_lo} set with \texttt{key\_hi}
still zero. This is sufficient on its own to cause root divergence
between validators.

These are both fixable inside 3.0 before 2025-12-25 with bounded local
changes. Critical \& high findings are all must-fix-pre-launch.

\begin{center}\rule{0.5\linewidth}{0.5pt}\end{center}

\subsection{Findings}\label{findings}

\subsubsection{Q3.0-STM-001 --- Commit ring backpressure produces
nondeterministic repair counts and divergent execution\_root
{[}Critical{]}}\label{q30-stm-001--commit-ring-backpressure-produces-nondeterministic-repair-counts-and-divergent-execution_root-critical}

\textbf{Description}. In \texttt{drain\_validate}, the order of
operations on a passing tx is:

\begin{enumerate}
\def\labelenumi{\arabic{enumi}.}
\tightlist
\item
  \texttt{mvcc\_apply\_writes(...)} (atomic version bump)
\item
  construct CommitItem, compute receipt\_hash
\item
  \texttt{ring\_try\_push(commit\_hdr,\ c)}
\item
  on push failure: requeue the \emph{original ExecResult} back onto
  Validate.
\end{enumerate}

The version bump in step 1 is permanent. If step 3 fails (commit ring
full --- explicitly called out as expected backpressure in the LP), the
tx will re-enter Validate on the next tick, but
\texttt{mvcc\_check\_consistent} will now find its own bumped version
(\texttt{cur\ !=\ e.version\_seen}) and route it to Repair. This
double-counts conflicts and repairs, AND alters which tx wins the keccak
chain in \texttt{drain\_commit}.

Worse: the comment in the code admits this is wrong:

\begin{quote}
"Roll back the MVCC version bump? In v0.36 substrate the version
monotonically advances; if commit ring is full we requeue and let
downstream catch up next tick. The version mismatch this introduces only
causes other concurrent txs to repair, which is correct under
Block-STM."
\end{quote}

It is \textbf{not} correct under Block-STM. It introduces a new conflict
that depends on whether the commit ring was full at that exact moment.
Two validators racing the GPU scheduler differently will fill the commit
ring at different points and produce different \texttt{repair\_count},
different commit ordering (different receipts/execution\_root), and
ultimately different \texttt{block\_hash}.

\textbf{Affected}: \texttt{quasar\_wave.metal:978-1004},
\texttt{quasar\_wave.cu:691-708}.

\textbf{Attack scenario}. Adversary submits a block where
validate-throughput exceeds commit-throughput for a sustained tick.
Easiest construction: a moderate flood of independent-key txs (no real
conflicts) + a tight \texttt{wave\_tick\_budget} so the commit drain
falls behind. On validator A the commit ring fills at tick K, on
validator B at tick K+1. Their receipts\_root chains observe a different
commit order. Roots diverge.

\textbf{Fix}. The commit-side push must be tried \emph{before} the MVCC
mutation, OR the mutation must be tentatively rolled back on push
failure. Cleanest fix:

\begin{enumerate}
\def\labelenumi{\arabic{enumi}.}
\tightlist
\item
  Try \texttt{ring\_try\_push(commit\_hdr,\ c)} first.
\item
  If it fails, push the ExecResult back to Validate without mutating
  MVCC. Tx will re-validate next tick and re-apply.
\end{enumerate}

Or: split commit into a "reserve commit slot" → "apply MVCC writes" →
"fill receipt" sequence that is atomic against the commit-ring-full
case.

\textbf{Must-fix pre-launch}: YES.

\begin{center}\rule{0.5\linewidth}{0.5pt}\end{center}

\subsubsection{\texorpdfstring{Q3.0-STM-002 --- Non-atomic key claim in
\texttt{mvcc\_locate} allows double-claim and torn key reads
{[}Critical{]}}{Q3.0-STM-002 --- Non-atomic key claim in mvcc\_locate allows double-claim and torn key reads {[}Critical{]}}}\label{q30-stm-002--non-atomic-key-claim-in-mvcc_locate-allows-double-claim-and-torn-key-reads-critical}

\textbf{Description}. \texttt{mvcc\_locate} claims an empty slot via
plain (non-atomic) stores:

\begin{verbatim}
if (s->key_lo == 0UL && s->key_hi == 0UL) {
    s->key_lo = key_lo;
    s->key_hi = key_hi;
    return idx;
}
\end{verbatim}

Two concurrent fibers landing on the same empty slot --- possible
because \texttt{drain\_exec} runs on workgroup \texttt{gid=4} but
\texttt{drain\_validate} (\texttt{gid=5}) and \texttt{drain\_dagready}
(\texttt{gid=3}) \textbf{also call \texttt{mvcc\_locate}} from their own
workgroups concurrently within one wave-tick dispatch. (Confirmed:
\texttt{drain\_dagready} doesn\textquotesingle t call
\texttt{mvcc\_locate}, but \texttt{drain\_exec},
\texttt{drain\_validate.mvcc\_check\_consistent}, and
\texttt{drain\_validate.mvcc\_apply\_writes} all do, from gid=4 and
gid=5 simultaneously. The kernel does NOT serialize these via
\texttt{gid==4} only.)

Even worse: the comment claims "races on (last\_writer\_tx, version)
atomics" make this safe, but those atomics only protect the version once
the slot is claimed. The \emph{claim itself} is racy:

\begin{itemize}
\tightlist
\item
  Thread T1 reads (0, 0), starts writing key\_lo = K.
\item
  Thread T2 reads (0, 0), also writes key\_lo = K, key\_hi = 0.
\item
  Thread T1 writes key\_hi = K\_HI.
\item
  Final state: key\_lo=K (correct), key\_hi=K\_HI (correct), version=0
  --- but T1 returned \texttt{idx} and T2 also returned \texttt{idx}.
  Both think they own the slot for \emph{different} keys if the keys
  happened to differ in the unlikely-collision case, OR for the same key
  (no harm) if they match.
\end{itemize}

The actually-exploitable case: \textbf{torn key read by another
concurrent fiber.} If a third thread T3 reaches the same idx during the
window between \texttt{s-\textgreater{}key\_lo\ =\ key\_lo} and
\texttt{s-\textgreater{}key\_hi\ =\ key\_hi}, T3 sees
\texttt{(key\_lo=K,\ key\_hi=0)} and either (a) misidentifies the slot
as "occupied by some other key (K,0)" and probes onward, or (b) finds
the comparator passes coincidentally and corrupts the wrong
key\textquotesingle s version.

\textbf{Affected}: \texttt{quasar\_wave.metal:489-510},
\texttt{quasar\_wave.cu:411-428}.

\textbf{Attack scenario}. Adversary submits a block whose first 16
transactions all touch \emph{different} keys whose FNV-1a hash collides
on the same MVCC slot. (FNV-1a is not cryptographic; finding 16
colliding keys is trivial --- pre-image / second-pre-image). Multiple
drain\_exec/validate workgroups race the empty slot, and depending on
GPU scheduling, the version chain on that slot is associated with a
nondeterministic key, producing different \texttt{state\_root} between
validators.

\textbf{Fix}. Use \texttt{atomic\_compare\_exchange} on a sentinel field
(e.g., on a packed (key\_lo\textbar busy) marker, or on a separate
\texttt{slot\_state} atomic) to claim the slot:

\begin{verbatim}
// Compete to install key_lo using atomic CAS from 0 → key_lo.
// On success, store key_hi (no race since only the winner reaches here),
// then publish the slot via a fence + store of a "ready" flag.
\end{verbatim}

Pair with \texttt{atomic\_thread\_fence(memory\_order\_acq\_rel)} to
publish the key bytes before any reader can match them.

Also: replace FNV-1a with a keyed hash (or use a structured key with
sufficient cryptographic mixing) so an adversary cannot trivially
hand-craft hash collisions. FNV-1a is fine for non-adversarial
workloads, not for a public chain accepting adversarial inputs.

\textbf{Must-fix pre-launch}: YES.

\begin{center}\rule{0.5\linewidth}{0.5pt}\end{center}

\subsubsection{\texorpdfstring{Q3.0-STM-003 --- Anti-livelock hard
limits (\texttt{MAX\_FAST\_REPAIRS}, \texttt{MAX\_TOTAL\_REPAIRS},
\texttt{HOT\_LANE\_CONFLICT\_THRESHOLD}) are not enforced anywhere in
code
{[}Critical{]}}{Q3.0-STM-003 --- Anti-livelock hard limits (MAX\_FAST\_REPAIRS, MAX\_TOTAL\_REPAIRS, HOT\_LANE\_CONFLICT\_THRESHOLD) are not enforced anywhere in code {[}Critical{]}}}\label{q30-stm-003--anti-livelock-hard-limits-max_fast_repairs-max_total_repairs-hot_lane_conflict_threshold-are-not-enforced-anywhere-in-code-critical}

\textbf{Description}. LP-010 §4.0 ("Anti-livelock hard limits ---
mandatory on GPU") documents:

\begin{verbatim}
constexpr uint32_t MAX_FAST_REPAIRS                = 3;
constexpr uint32_t MAX_TOTAL_REPAIRS               = 8;
constexpr uint32_t HOT_LANE_CONFLICT_THRESHOLD     = 16;
\end{verbatim}

with the caption: "Unbounded repair loops are unacceptable on GPU."

Grep across \texttt{quasar\_wave.metal}, \texttt{quasar\_wave.cu},
\texttt{quasar\_gpu\_engine.mm}, \texttt{quasar\_gpu\_engine\_cuda.cpp},
and \texttt{quasar\_gpu\_layout.hpp}: \textbf{none of these constants
exist}. There is no per-tx repair counter check, no
abort-on-repair-bound, no cold-queue demotion. \texttt{drain\_validate}
increments \texttt{result-\textgreater{}repair\_count} (a global
counter) and pushes the ExecResult back to Repair → Exec without any
bound.

The only thing limiting an N² repair storm is the host\textquotesingle s
choice of \texttt{max\_epochs} in \texttt{run\_until\_done}. The test
\texttt{test\_block\_stm\_conflict\_repair} passes
\texttt{max\_epochs=256} and asserts only that
\texttt{r.status\ ==\ 1u}. If 256 ticks are insufficient (because the
adversary submits a block that requires more), the round fails
non-deterministically: validator A might have a fast GPU and finish in
200 ticks; validator B with a slower GPU times out at 256.

\textbf{Affected}: entire kernel --- there is nowhere to put the counter
check because it doesn\textquotesingle t exist.

\textbf{Attack scenario}. Adversary submits 1024 transactions all
writing the same key. Theoretical bound is O(N²) ≈ 524K repairs; even
with an aggressive 1M ops/sec drain that\textquotesingle s 0.5 seconds
of pure repair churn --- but the worst-case bound is actually higher
because repair → exec → validate can re-enter validation without the
original tx ever committing (the schedule is GPU-nondeterministic).
Easier attack: 256 same-key txs guaranteed to produce \textasciitilde32K
repairs. The adversary picks block size to land just past whatever the
host chose for \texttt{max\_epochs}. Some validators fail to finalize;
others succeed; chain splits.

\textbf{Fix}. Add \texttt{repair\_count} and \texttt{incarnation} to the
FiberSlot or a per-tx state record. In \texttt{drain\_validate}, before
pushing to Repair:

\begin{verbatim}
if (er.incarnation >= MAX_TOTAL_REPAIRS) {
    // emit deterministic abort: status = 4 (Error), zero gas, deterministic receipt
    // this is the "deterministic block-pathological-tx" outcome
}
\end{verbatim}

The abort path must be deterministic --- same input must produce the
same abort across validators --- which means the threshold must be
enforced \emph{and} the abort must produce a fixed receipt regardless of
how many wave-ticks were consumed.

Also publish \texttt{MAX\_TOTAL\_REPAIRS} to consensus, because changing
it is a hard fork.

\textbf{Must-fix pre-launch}: YES.

\begin{center}\rule{0.5\linewidth}{0.5pt}\end{center}

\subsubsection{Q3.0-STM-004 --- Cross-backend determinism is unverified;
CUDA backend lacks v0.40 DAG construction
{[}Critical{]}}\label{q30-stm-004--cross-backend-determinism-is-unverified-cuda-backend-lacks-v040-dag-construction-critical}

\textbf{Description}. LP-010 §"Correctness Invariants" \#7 states:

\begin{quote}
Cross-backend equivalence --- Metal == CUDA == CPU reference for
state\_root, receipts\_root, execution\_root, mode\_root, and
conflict/repair behavior wherever deterministic.
\end{quote}

Reality:

\begin{enumerate}
\def\labelenumi{\arabic{enumi}.}
\tightlist
\item
  The CUDA backend (\texttt{quasar\_wave.cu:571-587}) ships \textbf{only
  the v0.35 pass-through \texttt{drain\_dagready}}. It has no
  DagWriterSlot table, no PredictedKey lookup, no parent→child edge
  construction. The Metal backend has full v0.40 DAG construction
  (\texttt{quasar\_wave.metal:736-825}).
\item
  The CUDA \texttt{drain\_commit} (\texttt{quasar\_wave.cu:770-809})
  does NOT walk DAG children to decrement unresolved\_parents. The Metal
  \texttt{drain\_commit} (\texttt{quasar\_wave.metal:1107-1127}) does.
\item
  In Nebula mode (\texttt{desc-\textgreater{}mode\ ==\ 1u}), Metal
  builds a real DAG and uses unresolved\_parents to gate emission to
  Exec; CUDA pass-through\textquotesingle s antichain is just FIFO.
\item
  The CUDA vote verifier (\texttt{drain\_vote} in
  \texttt{quasar\_wave.cu:811-860}) uses a different signature scheme
  entirely (\texttt{verify\_signature\_stub} byte-equality check on
  subject vs sig) than the Metal vote-batch verifier (HMAC-keccak with
  master secret + chain\_id binding).
\end{enumerate}

\textbf{Two validators, one on Metal (Apple node) and one on CUDA (Linux
x86), running the same Nebula-mode block, will produce different
execution\_roots} because the order in which children are released to
Exec is different.

There is also no cross-backend test harness: no
\texttt{cevm/test/cross\_backend/...}, no CI step that runs the same
block through both kernels and diff\textquotesingle d the roots.

\textbf{Affected}: \texttt{quasar\_wave.cu:571-587},
\texttt{quasar\_wave.cu:770-809}, \texttt{quasar\_wave.cu:811-860}, all
of vote verification, missing test harness.

\textbf{Attack scenario}. Geographically separated validators have
heterogeneous GPUs (Apple silicon vs NVIDIA --- explicit goal of the
multi-backend design). Adversary submits a Nebula-mode block with a
multi-key DAG. Apple-side validators emit DAG children in the order
computed by drain\_commit\textquotesingle s child walk; CUDA-side
validators emit them in FIFO order from the dagready ring. Different
commit order → different keccak chain → different roots → chain split.

\textbf{Fix}.

\begin{enumerate}
\def\labelenumi{\arabic{enumi}.}
\tightlist
\item
  Port v0.40 drain\_dagready and drain\_commit-with-children to CUDA.
  Add the same DagWriterSlot table and the same PredictedKey lookup.
  Same atomic and same fence semantics.
\item
  Port the HMAC-keccak vote verifier (matching kernel signatures must
  match across backends --- the byte-equality stub is also a critical
  security bug, see Q3.0-STM-008).
\item
  Build a CPU reference (LP-010 v0.49 promises this --- must land for
  2025-12-25).
\item
  Add \texttt{cevm/test/cross\_backend/quasar\_round\_diff\_test.cpp}
  that runs the same QuasarRoundDescriptor + tx batch through Metal,
  CUDA, and CPU reference, asserts byte-identical roots.
\item
  CI must reject any commit that produces different roots across
  backends.
\end{enumerate}

\textbf{Must-fix pre-launch}: YES --- this is the difference between
"GPU-accelerated chain" and "permanent fork on heterogeneous hardware."

\begin{center}\rule{0.5\linewidth}{0.5pt}\end{center}

\subsubsection{\texorpdfstring{Q3.0-STM-005 --- \texttt{drain\_dagready}
cross-workgroup race on \texttt{last\_writer\_tx} and
\texttt{unresolved\_parents} allows missed wake-ups
{[}High{]}}{Q3.0-STM-005 --- drain\_dagready cross-workgroup race on last\_writer\_tx and unresolved\_parents allows missed wake-ups {[}High{]}}}\label{q30-stm-005--drain_dagready-cross-workgroup-race-on-last_writer_tx-and-unresolved_parents-allows-missed-wake-ups-high}

\textbf{Description}. The Metal kernel comments above
\texttt{DagWriterSlot} claim "drain\_dagready is single-threaded (gid=3,
tid=0), so last\_writer\_tx is plain (not atomic)" --- but the protocol
still has a real race against \texttt{drain\_commit} (gid=7, tid=0)
which concurrently:

\begin{itemize}
\tightlist
\item
  Walks \texttt{dag\_nodes{[}c.tx\_index{]}.children} (line 1108-1124)
\item
  Calls
  \texttt{atomic\_fetch\_sub\_explicit(\&C-\textgreater{}unresolved\_parents,\ 1u)}
  and re-emits child to Exec
\item
  Stores \texttt{state\ =\ kDagNodeCommitted}
\end{itemize}

drain\_dagready (gid=3) reads \texttt{dag\_nodes{[}prev{]}.state} (line
788), increments \texttt{T-\textgreater{}unresolved\_parents} (line
790), and \emph{only then} checks state again at line 797. The
double-check is racy:

\begin{verbatim}
1. drain_dagready reads P->state == Registered          (T1 at gid=3)
2. drain_commit transitions P → Committed and walks
   children[*]; the walk completes BEFORE drain_dagready
   has incremented T->unresolved_parents               (T2 at gid=7)
3. drain_dagready increments T->unresolved_parents      (T1)
4. drain_dagready reads P->state == Committed and 
   decrements T->unresolved_parents                     (T1)
5. T->unresolved_parents settles at 0 — correct in
   isolation, but T's state never transitioned to 
   Emitted because drain_dagready's emit-when-zero 
   check at line 814 fires only on the increment path.
\end{verbatim}

Specifically: at line 815, \texttt{unresolved} is read after the
for-loop processing all parents. If the very last
parent\textquotesingle s \texttt{state} flipped to Committed
\emph{between} line 798\textquotesingle s check and line
814\textquotesingle s read, but drain\_commit\textquotesingle s
child-walk of that last parent had already completed (it walked children
before transitioning to Committed at line 1126, but the walk is
per-parent, not per-(parent, predicted-edge-key)), then
drain\_commit\textquotesingle s \texttt{atomic\_fetch\_sub\_explicit}
ran on a \texttt{unresolved\_parents} value that did NOT include this
child\textquotesingle s increment yet, and the "old == 1u" check at line
1118 fires on an old value that\textquotesingle s wrong. Children can
either (a) miss being re-emitted, or (b) be double-emitted (race between
drain\_dagready\textquotesingle s emit-when-zero at line 814 and
drain\_commit\textquotesingle s emit-when-zero at line 1120).

\textbf{Affected}: \texttt{quasar\_wave.metal:736-825}
(drain\_dagready), \texttt{quasar\_wave.metal:1107-1127} (drain\_commit
DAG child walk).

\textbf{Attack scenario}. Adversary submits a Nebula block with
carefully timed parent/child relationships such that
drain\_commit\textquotesingle s walk of P\textquotesingle s children
completes exactly when drain\_dagready is processing a new child of P.
Validator A\textquotesingle s GPU schedules drain\_commit slightly
earlier; validator B\textquotesingle s slightly later. Validator A
re-emits the child once (correct). Validator B re-emits twice OR not at
all. tx\_count diverges.

\textbf{Fix}. The drain\_commit child walk and the drain\_dagready
edge-add must be linearized. Options:

\begin{enumerate}
\def\labelenumi{\arabic{enumi}.}
\tightlist
\item
  Move drain\_commit\textquotesingle s child-walk and state-transition
  into a CAS-based protocol: child increments unresolved, then
  drain\_commit\textquotesingle s decrement+emit only fires if it
  observes the increment. The current "double-check after store" pattern
  in drain\_dagready is correct \emph{if} drain\_commit uses release
  semantics around \texttt{state\ =\ Committed} and a
  \texttt{mem\_device} fence between child-walk and state-store. The
  kernel currently uses \texttt{memory\_order\_relaxed} for both state
  load and store. \textbf{That is the bug.}
\item
  Use \texttt{memory\_order\_release} for \texttt{state\ =\ Committed}
  after the child-walk, and \texttt{memory\_order\_acquire} for the
  second \texttt{state} read in drain\_dagready. The C++ memory model
  guarantees the new edge --- added before the increment --- is visible
  to drain\_commit.
\end{enumerate}

In Metal, replace \texttt{memory\_order\_relaxed} with
\texttt{memory\_order\_release}/\texttt{memory\_order\_acquire} on the
\texttt{state} atomic. Same for CUDA\textquotesingle s
\texttt{\_\_threadfence()} placement.

\textbf{Must-fix pre-launch}: YES (high). Determinism violation under
Nebula mode.

\begin{center}\rule{0.5\linewidth}{0.5pt}\end{center}

\subsubsection{\texorpdfstring{Q3.0-STM-006 ---
\texttt{predicted\_keys{[}tx\_index{]}} indexing assumes drain\_ingress
assigns sequential tx\_index, but the host pre-fills before GPU consumes
{[}High{]}}{Q3.0-STM-006 --- predicted\_keys{[}tx\_index{]} indexing assumes drain\_ingress assigns sequential tx\_index, but the host pre-fills before GPU consumes {[}High{]}}}\label{q30-stm-006--predicted_keystx_index-indexing-assumes-drain_ingress-assigns-sequential-tx_index-but-the-host-pre-fills-before-gpu-consumes-high}

\textbf{Description}. The host driver
(\texttt{quasar\_gpu\_engine.mm:276-291}) writes:

\begin{Shaded}
\begin{Highlighting}[]
\AttributeTok{const} \DataTypeTok{uint32\_t}\NormalTok{ tidx }\OperatorTok{=} \VariableTok{round\_}\OperatorTok{.}\NormalTok{next\_predicted\_idx}\OperatorTok{;}
\ControlFlowTok{if} \OperatorTok{(}\NormalTok{tidx }\OperatorTok{\textless{}}\NormalTok{ kDagNodeCapacity}\OperatorTok{)} \OperatorTok{\{}
\NormalTok{    PredictedKeyHost}\OperatorTok{*}\NormalTok{ slot }\OperatorTok{=} \OperatorTok{\&}\NormalTok{predicted\_arena}\OperatorTok{[}\NormalTok{tidx }\OperatorTok{*}\NormalTok{ kMaxPredictedKeys}\OperatorTok{];}
    \OperatorTok{...}
\OperatorTok{\}}
\VariableTok{round\_}\OperatorTok{.}\NormalTok{next\_predicted\_idx }\OperatorTok{+=} \DecValTok{1}\BuiltInTok{u}\OperatorTok{;}
\end{Highlighting}
\end{Shaded}

This assumes \texttt{tx\_index\_seq} on the GPU side will exactly match
\texttt{next\_predicted\_idx} on the host side. The GPU assigns
\texttt{tx\_index\ =\ atomic\_fetch\_add(tx\_index\_seq,\ 1)} in
drain\_ingress (\texttt{quasar\_wave.metal:527}). drain\_ingress runs
single-threaded on \texttt{gid=0}, FIFO over Ingress, so the assignment
IS deterministic FIFO --- but only as long as:

\begin{enumerate}
\def\labelenumi{\arabic{enumi}.}
\tightlist
\item
  The host pushes txs in the same order as it pre-fills
  \texttt{predicted\_keys}.
\item
  No tx is dropped at admission (e.g., decode failure routes back to
  \texttt{(void)ring\_try\_push(ingress\_hdr,\ ...)} on overflow).
\item
  Across multiple \texttt{push\_txs(...)} calls, no other writer
  pre-fills predicted\_keys at indices the host hasn\textquotesingle t
  reserved yet.
\end{enumerate}

Concrete failure: on lines 535-538 of drain\_ingress, if
\texttt{ring\_try\_push(decode\_hdr,\ ...)} fails (decode ring full),
the pop\textquotesingle d tx is requeued to Ingress with
\texttt{(void)ring\_try\_push(ingress\_hdr,\ ingress\_items,\ in)}.
\textbf{But the \texttt{tx\_index} was already consumed by
\texttt{atomic\_fetch\_add(tx\_index\_seq,\ 1)} at line 527.} On the
next tick that same tx pops back, gets a \emph{new} tx\_index, and the
slot at the old tx\_index is \texttt{predicted\_keys} indexed by the old
(orphaned) value. The slot at the \emph{new} tx\_index belongs to a
different tx that hasn\textquotesingle t been pushed yet --- so
\texttt{predicted\_keys{[}new\_tx\_index{]}} reads zeroed data. The DAG
misses the predicted dependency.

Also: if \texttt{push\_txs} is called multiple times within a round,
\texttt{next\_predicted\_idx} increments host-side per tx pushed, but if
the \emph{first} push\textquotesingle s drain\_ingress fails the decode
push (some txs successfully decoded, some requeued), the host side
believes tx\_index=0..N-1 has the predicted\_access for txs 0..N-1, but
the GPU\textquotesingle s tx\_index\_seq is at 0..N-K (some requeued,
will get index N..N+K-1 next round). The mapping is broken.

\textbf{Affected}: \texttt{quasar\_gpu\_engine.mm:276-291},
\texttt{quasar\_wave.metal:527,\ 535-538}.

\textbf{Attack scenario}. Adversary spam-fills the decode ring in early
ticks so drain\_ingress requeues txs. The orphaned tx\_indexes leave
gaps in tx\_index space; predicted\_access lookup misses real edges; DAG
construction omits dependencies; downstream Block-STM detects the
conflict and forces repairs OR worse, commits in the wrong order.
Different validators hit the requeue at different ticks → different
tx\_index assignments → different DAG → different roots.

\textbf{Fix}. Either:

\begin{enumerate}
\def\labelenumi{\arabic{enumi}.}
\tightlist
\item
  Move \texttt{tx\_index\_seq} increment to \emph{after} successful
  decode push, not before. (Cleanest.)
\item
  Index predicted\_keys by a stable key (origin \textbar\textbar{}
  nonce) rather than tx\_index.
\item
  Use a per-round content-addressed predicted\_access lookup (hash of tx
  envelope → predicted\_keys offset).
\end{enumerate}

Option 1 is the smallest change.

\textbf{Must-fix pre-launch}: YES (high). Causes Nebula determinism
violation under any backpressure.

\begin{center}\rule{0.5\linewidth}{0.5pt}\end{center}

\subsubsection{\texorpdfstring{Q3.0-STM-007 --- \texttt{drain\_vote}
reads \texttt{vote\_verified{[}head\_pre\ \&\ mask{]}} AFTER popping;
race with verifier produces accept-tampered-vote
{[}High{]}}{Q3.0-STM-007 --- drain\_vote reads vote\_verified{[}head\_pre \& mask{]} AFTER popping; race with verifier produces accept-tampered-vote {[}High{]}}}\label{q30-stm-007--drain_vote-reads-vote_verifiedhead_pre--mask-after-popping-race-with-verifier-produces-accept-tampered-vote-high}

\textbf{Description}. The Metal \texttt{drain\_vote} (line 1148-1187)
reads:

\begin{verbatim}
uint head_pre = atomic_load_explicit(&vote_hdr->head, memory_order_relaxed);
VoteIngress v;
if (!ring_try_pop(vote_hdr, vote_items, v)) break;
uint vidx = head_pre & (vote_verified_capacity - 1u);
if (vote_verified[vidx] == 0u) { ++processed; continue; }
\end{verbatim}

Two issues:

\begin{enumerate}
\def\labelenumi{\arabic{enumi}.}
\tightlist
\item
  \texttt{head\_pre} is read \emph{before} \texttt{ring\_try\_pop}
  advances head. ring\_try\_pop loops on CAS --- if multiple drain\_vote
  workgroups raced, the \texttt{head\_pre} here might NOT match the slot
  that was actually popped. (In current dispatch, \texttt{gid=10} runs
  single-threaded with \texttt{tid==0}, so this works in v0.38 --- but
  it is fragile and assumes the dispatcher never schedules two threads
  on gid=10.)
\item
  The vote\_verified bitmap is written by
  \texttt{quasar\_verify\_votes\_kernel} BEFORE the wave-tick scheduler
  runs (line 437-448 in \texttt{quasar\_gpu\_engine.mm} --- separate
  encoder, ordered within one MTLCommandBuffer). The verifier writes
  verified{[}slot\_idx{]} = 1 if signature matches. \textbf{But the
  wave-tick scheduler also runs \texttt{drain\_ingress} which advances
  tx\_index\_seq, and inside the same dispatch can call
  \texttt{drain\_vote} against a vote\_verified bitmap computed for the
  previous tail --- meaning a vote pushed \emph{during this run\_epoch}
  is popped without ever being verified.} Specifically: host pushes vote
  V at tail=T. The verifier ran at tail\textless T, so
  vote\_verified{[}T mask{]} is stale (might be 0 from initialization,
  OR might be 1 from a stale match for a previous vote at the same
  slot).
\end{enumerate}

Combined, the bitmap can return either:

\begin{itemize}
\tightlist
\item
  0 → vote rejected (best-case false negative, denies vote even though
  sig is valid)
\item
  1 → vote accepted (worst-case false positive: stale verified=1 from a
  previous round/vote at the same slot causes a tampered/replayed vote
  to count for stake)
\end{itemize}

\textbf{Affected}: \texttt{quasar\_wave.metal:1148-1154},
\texttt{quasar\_gpu\_engine.mm:430-475}.

\textbf{Attack scenario}. Adversary controls a validator account with
low stake. Adversary submits two consecutive rounds. Round N: legitimate
vote at slot S, verified=1. Round N+1: adversary submits a forged vote
(different validator, different signature) at slot S\textquotesingle{}
where S\textquotesingle{} \& mask == S \& mask. If the host
doesn\textquotesingle t memset vote\_verified between rounds
(\texttt{begin\_round} does memset at line 212, so this specific case is
safe per-round), or if the same round has high vote churn (legitimate
vote at slot S; verifier marks verified{[}S{]}=1; vote consumed;
\emph{new} vote pushed at slot S+capacity which masks to S; verifier
hasn\textquotesingle t run yet so verified{[}S{]}=1 from before but now
refers to a different vote), the forged vote is accepted. Stake count
diverges between validators with different scheduling. Quorum
certificate emission depends on stake counts. Two validators emit
different certs. Worse: a malicious validator can use this to inflate
own stake via cross-round bitmap pollution.

\textbf{Affected}: \texttt{quasar\_wave.metal:1148-1187}.

\textbf{Fix}. The vote-verified bitmap must be (a) sized larger than the
ring (to avoid mask collisions across in-flight votes), (b) cleared
between verifier runs OR the verifier must run after every push, (c)
keyed by a content hash of the vote rather than slot index, and (d) the
verifier and consumer must be ordered against any new vote arrivals via
the same ring\textquotesingle s pushed counter.

Cleanest: re-run the verifier inline as part of drain\_vote --- verify
the signature there and skip the bitmap entirely. The bitmap is a
premature optimization that introduces a TOCTOU.

\textbf{Must-fix pre-launch}: YES --- quorum integrity is a chain-safety
primitive.

\begin{center}\rule{0.5\linewidth}{0.5pt}\end{center}

\subsubsection{\texorpdfstring{Q3.0-STM-008 --- CUDA
\texttt{verify\_signature\_stub} is byte-equality (sig == subject);
accepts any signature where signature{[}0..32{]} == subject
{[}High{]}}{Q3.0-STM-008 --- CUDA verify\_signature\_stub is byte-equality (sig == subject); accepts any signature where signature{[}0..32{]} == subject {[}High{]}}}\label{q30-stm-008--cuda-verify_signature_stub-is-byte-equality-sig--subject-accepts-any-signature-where-signature032--subject-high}

\textbf{Description}. CUDA \texttt{quasar\_wave.cu:811-817}:

\begin{verbatim}
__device__ __forceinline__ bool verify_signature_stub(const uint8_t* subject, const uint8_t* sig)
{
    for (uint32_t i = 0u; i < 32u; ++i) {
        if (sig[i] != subject[i]) return false;
    }
    return true;
}
\end{verbatim}

This accepts any vote where
\texttt{signature{[}0..32{]}\ ==\ subject{[}0..32{]}}. An attacker who
knows the subject (it\textquotesingle s the block\_hash, broadcast to
all validators) can construct a "signature" by literally copying subject
into the first 32 bytes of the signature field. Result: anyone can forge
votes for any validator on the CUDA backend.

The Metal backend uses \texttt{quasar\_verify\_votes\_kernel} with
HMAC-keccak against a \texttt{kQuasarMasterSecret} --- also broken (the
master secret is hardcoded in the kernel source and is therefore public,
see Q3.0-STM-009), but at least it\textquotesingle s not a one-line
equality check.

\textbf{Affected}: \texttt{quasar\_wave.cu:811-817}.

\textbf{Attack scenario}. CUDA validator. Adversary observes block\_hash
via the network (this is broadcast). Constructs a vote message with
\texttt{signature{[}0..32{]}\ =\ block\_hash{[}0..32{]}}, any garbage in
\texttt{{[}32..96{]}}, any validator\_index, any stake\_weight.
drain\_vote accepts the vote, accumulates stake\_weight, emits a
QuorumCert for that lane at quorum threshold. \textbf{Adversary controls
all three lanes (BLS / Ringtail / MLDSA) trivially.}

\textbf{Fix}. The CUDA verifier must implement the same HMAC-keccak path
as Metal (still broken, see next finding). Long-term: real BLS12-381 /
Ringtail / ML-DSA verifiers must be CUDA-ported too.

\textbf{Must-fix pre-launch}: YES --- this is a one-line attack to forge
consensus signatures on CUDA validators.

\begin{center}\rule{0.5\linewidth}{0.5pt}\end{center}

\subsubsection{\texorpdfstring{Q3.0-STM-009 ---
\texttt{kQuasarMasterSecret} is a hardcoded constant in the kernel
source; HMAC-keccak verifier provides zero security
{[}High{]}}{Q3.0-STM-009 --- kQuasarMasterSecret is a hardcoded constant in the kernel source; HMAC-keccak verifier provides zero security {[}High{]}}}\label{q30-stm-009--kquasarmastersecret-is-a-hardcoded-constant-in-the-kernel-source-hmac-keccak-verifier-provides-zero-security-high}

\textbf{Description}. \texttt{quasar\_wave.metal:1212-1215}:

\begin{verbatim}
constant uchar kQuasarMasterSecret[32] = {
    0x51,0x55,0x41,0x53,0x41,0x52,0x2D,0x76,0x30,0x33,0x38,0x2D,0x6D,0x61,0x73,0x74,
    0x65,0x72,0x2D,0x73,0x65,0x63,0x72,0x65,0x74,0x2D,0x73,0x68,0x61,0x72,0x65,0x64,
};
\end{verbatim}

ASCII: \texttt{"QUASAR-v038-master-secret-shared"}.

This constant is in the kernel source, the LP, and presumably the
binary. The "verification"
\texttt{verified{[}slot\_idx{]}\ =\ (HMAC(kQuasarMasterSecret,\ validator\_index\textbar{}\textbar{}subject\textbar{}\textbar{}round)\ ==\ signature)}
is therefore:

\begin{quote}
"anyone with access to the kernel source can compute the expected
signature for any (validator, subject, round)."
\end{quote}

The kernel source is open source (Apache-2.0 per the file header).
\textbf{The signature verifier rejects honest validators who use real
BLS signatures and accepts attackers who compute the HMAC-keccak.}

The LP-132 documentation acknowledges this is a placeholder for v0.43+,
but the launch is 2025-12-25 and the LP\textquotesingle s roadmap shows
real BLS at "v0.43" with no commitment to real Ringtail (v0.44) or
Groth16 (v0.45) by launch.

\textbf{Affected}: \texttt{quasar\_wave.metal:1212-1290},
\texttt{quasar\_gpu\_engine.mm:430-448} (verify\_pso\_),
\texttt{quasar\_wave.cu:811-860}.

\textbf{Attack scenario}. Adversary clones the repo, reads the master
secret, computes valid signatures for every validator, posts votes for
every lane → controls all quorums. Cost: zero.

\textbf{Fix}. The launch MUST ship real BLS12-381 pairing for the BLS
lane (LP-075), real Ringtail share verification for the RT lane
(LP-073), and real ML-DSA-65 (LP-070) for the MLDSA lane. Anything less
is a ceremonial signature. If real BLS is not ready by 2025-12-25, the
launch must NOT enable on-chain quorum aggregation in QuasarGPU --- fall
back to host-side verification using vetted libraries
(\texttt{luxfi/crypto}).

\textbf{Must-fix pre-launch}: YES --- without this, consensus is
theater.

\begin{center}\rule{0.5\linewidth}{0.5pt}\end{center}

\subsubsection{\texorpdfstring{Q3.0-STM-010 --- Closing-gate
\texttt{commit\_consumed\ ==\ ingress\_pushed} doesn\textquotesingle t
account for txs in-flight in StateRequest, Decode, Crypto, etc.
{[}Medium{]}}{Q3.0-STM-010 --- Closing-gate commit\_consumed == ingress\_pushed doesn\textquotesingle t account for txs in-flight in StateRequest, Decode, Crypto, etc. {[}Medium{]}}}\label{q30-stm-010--closing-gate-commit_consumed--ingress_pushed-doesnt-account-for-txs-in-flight-in-staterequest-decode-crypto-etc-medium}

\textbf{Description}. The finalization gate
(\texttt{quasar\_wave.metal:1389-1414}):

\begin{verbatim}
if (gid == 0u && desc->closing_flag != 0u) {
    const uint ingress_pushed = atomic_load_explicit(&hdrs[0].pushed);
    const uint commit_consumed = atomic_load_explicit(&hdrs[7].consumed);
    if (ingress_pushed == commit_consumed) {
        // ... compute block_hash, set status=1 (finalized) ...
    }
}
\end{verbatim}

The gate compares \texttt{ingress.pushed} against
\texttt{commit.consumed}. But a tx that goes Ingress → Decode →
StateRequest (cold-state) and is awaiting a host StatePage will have:

\begin{itemize}
\tightlist
\item
  ingress.pushed = N (incremented on push)
\item
  commit.consumed = N-1 (this tx hasn\textquotesingle t committed)
\end{itemize}

This is the intended behavior --- the round won\textquotesingle t
finalize until all txs commit. Good.

However: \texttt{drain\_decode} increments \texttt{fibers\_suspended}
when routing a tx to StateRequest, and \texttt{drain\_state\_resp}
re-injects via \texttt{crypto\_hdr} (line 1057) --- where it re-enters
\texttt{drain\_crypto}. \textbf{drain\_crypto increments
\texttt{processed} but NOT \texttt{pushed} on Ingress.} So
ingress.pushed stays at N (correct --- the tx was originally ingressed).
But what if \texttt{push\_txs} is called again \emph{during} a round to
add a new tx after closing\_flag was set? The host drives
\texttt{request\_close()} setting closing\_flag, but \texttt{push\_txs}
is not gated against closing --- line 242-298 of
\texttt{quasar\_gpu\_engine.mm} doesn\textquotesingle t check
\texttt{closing\_flag}. The host can push txs after setting
closing\_flag, ingress.pushed grows, the gate doesn\textquotesingle t
fire on the next tick (correct), but tx ordering is now nondeterministic
if two validators differ in when their host pushes the late tx.

\textbf{Affected}: \texttt{quasar\_gpu\_engine.mm:242-298} (push\_txs
missing closing\_flag guard), \texttt{quasar\_wave.metal:1389-1414}
(gate).

\textbf{Attack scenario}. Two validators receive the same block.
Validator A\textquotesingle s \texttt{request\_close()} and
\texttt{push\_txs()} happen in order A; validator B\textquotesingle s
happen in order B with race. The late-pushed tx ends up in different
positions in the canonical order across validators. Roots diverge.

\textbf{Fix}. \texttt{push\_txs} must reject pushes after
\texttt{closing\_flag} is set, returning a deterministic error to the
host. The host must validate tx\_count against the round
descriptor\textquotesingle s expected count before calling
\texttt{request\_close()}.

\textbf{Must-fix pre-launch}: YES (medium-rising-to-high). Fix is two
lines.

\begin{center}\rule{0.5\linewidth}{0.5pt}\end{center}

\subsubsection{\texorpdfstring{Q3.0-STM-011 --- \texttt{state\_root} is
never written by any drain; \texttt{block\_hash} includes uninitialized
state\_root
{[}Medium{]}}{Q3.0-STM-011 --- state\_root is never written by any drain; block\_hash includes uninitialized state\_root {[}Medium{]}}}\label{q30-stm-011--state_root-is-never-written-by-any-drain-block_hash-includes-uninitialized-state_root-medium}

\textbf{Description}. The finalization gate computes:

\begin{verbatim}
block_hash = keccak(round || mode || receipts_root || execution_root || state_root || mode_root)
\end{verbatim}

But \texttt{state\_root} is never written by any drain. The kernel uses
\texttt{result-\textgreater{}state\_root{[}k{]}} at line 1404 ---
read-only. Grep for \texttt{state\_root{[}} writes: only the line 1404
read.

\texttt{begin\_round} \texttt{memset}s \texttt{result\_buf} to zero at
line 205, so state\_root is all-zero on every round. The block\_hash on
every round includes the constant 32-zero-bytes for state\_root,
regardless of what state was actually mutated. This is deterministic
across runs (good for the determinism test), but it \textbf{provides no
commitment to the actual MVCC state}. Two distinct executions producing
the same receipts/execution traces but different actual state commit to
the same block\_hash.

This is the definition of a "ceremony hash" --- the block\_hash binds
nothing about state.

\textbf{Affected}: \texttt{quasar\_wave.metal:1404}, all drains (no
writer).

\textbf{Attack scenario}. Adversary submits block X. Honest validators
execute it, producing some MVCC state. Adversary then constructs a block
Y that produces the same receipts (same external observations) but
different MVCC state. \textbf{Both produce the same block\_hash.} State
machines fork silently on internal state; later txs depending on state
diverge between A and B, but block\_hash equality means the consensus
engine sees no fork.

\textbf{Fix}. After the validate-commit pipeline drains, run a
\texttt{drain\_state\_root} workgroup (or extend drain\_commit) that
computes:

\begin{verbatim}
state_root = keccak(running_state_root || (key_lo || key_hi || version || last_writer_tx)_in_canonical_order)
\end{verbatim}

The canonical order MUST be by \texttt{(tx\_index,\ incarnation)} of the
committing tx --- which is exactly what\textquotesingle s available in
CommitItem. Or: maintain a Merkle-style accumulator over MVCC slot
deltas as drain\_commit fires.

Note that the LP-010 §"Three-Tier Validation" and §"MVCC Garbage
Collection" promise this --- it\textquotesingle s just not implemented.

\textbf{Must-fix pre-launch}: YES (medium). The block\_hash without a
real state\_root is not a usable consensus commitment.

\begin{center}\rule{0.5\linewidth}{0.5pt}\end{center}

\subsubsection{Q3.0-STM-012 --- Cold-state suspend/resume race:
re-injected fiber may run with stale "loaded sentinel"
{[}Medium{]}}\label{q30-stm-012--cold-state-suspendresume-race-re-injected-fiber-may-run-with-stale-loaded-sentinel-medium}

\textbf{Description}. LP-132 §"Async cold-state page faults" describes
the loaded sentinel:

\begin{quote}
drain\_state\_resp resumes by re-injecting and stamping the MVCC
slot\textquotesingle s
\texttt{last\_writer\_tx\ \textbar{}=\ 0x80000000} as the "loaded"
sentinel.
\end{quote}

In code: \texttt{drain\_state\_resp} at
\texttt{quasar\_wave.metal:1041-1065} re-injects via crypto\_hdr but
does NOT touch the MVCC slot\textquotesingle s
\texttt{last\_writer\_tx}. The "loaded sentinel" stamping is documented
in LP-132 but not implemented. Without it:

\begin{enumerate}
\def\labelenumi{\arabic{enumi}.}
\tightlist
\item
  tx T1 reads cold key K, suspends → StateRequest emitted.
\item
  Host services StateRequest, posts StatePage.
\item
  drain\_state\_resp re-injects T1 via crypto\_hdr → drain\_crypto →
  drain\_exec.
\item
  drain\_exec calls \texttt{mvcc\_locate(K)} which finds the slot empty
  (because no tx has written K yet on the GPU side), claims it via the
  racy logic in Q3.0-STM-002, reads \texttt{version=0}, builds RW set
  with \texttt{version\_seen=0}.
\item
  tx T2 with same K already executed in steps 2-3 (because drain\_exec
  doesn\textquotesingle t block on T1), wrote version=1.
\item
  T1\textquotesingle s validate fails (cur=1, version\_seen=0), repair,
  re-exec. \textbf{But T1\textquotesingle s host page data is gone ---
  it was a one-shot StatePage.}
\end{enumerate}

The kernel will spin on T1: every repair re-exec finds version mismatch,
T1 never sees the actual loaded data because the StatePage payload
isn\textquotesingle t preserved across repairs. Without the loaded
sentinel pinning T1\textquotesingle s first execution to commit before
T2, T1 starves.

Worse: two concurrent host StatePages for the same tx\_index (e.g.,
kernel emitted two StateRequests for two different cold keys for the
same tx) --- the kernel does not coalesce them. drain\_state\_resp
re-injects on every page received, so a tx that needs two cold pages
will be re-injected twice into crypto\_hdr → executed twice →
double-counted in drain\_commit → tx\_count over by one.

\textbf{Affected}: \texttt{quasar\_wave.metal:1041-1065}
(drain\_state\_resp), \texttt{quasar\_wave.metal:891-945} (drain\_exec
--- no per-tx pending state machine),
\texttt{quasar\_gpu\_engine.mm:359-390} (push\_state\_pages allows
multiple per tx).

\textbf{Attack scenario}. Adversary submits txs that touch multiple cold
keys (force multiple StateRequests per tx) AND target hot lanes (for
repairs). Validator A\textquotesingle s host serves pages in order P1,
P2; validator B in order P2, P1. Different re-injection orders →
different tx\_index ordering at validate → different roots.

\textbf{Fix}.

\begin{enumerate}
\def\labelenumi{\arabic{enumi}.}
\tightlist
\item
  Implement the loaded sentinel: in drain\_state\_resp, before
  re-injecting, stamp
  \texttt{mvcc\_table{[}K{]}.last\_writer\_tx\ \textbar{}=\ 0x80000000}
  (or use a separate per-slot loaded flag with the high bit reserved).
  drain\_validate must respect this sentinel and pin the resumed
  tx\textquotesingle s commit to happen before any other writer of K.
\item
  Coalesce multiple StatePage arrivals per tx\_index into a single
  re-injection. Track per-tx pending\_count in FiberSlot; only re-inject
  when pending\_count reaches zero.
\item
  Clear the loaded sentinel on commit, never on repair.
\end{enumerate}

\textbf{Must-fix pre-launch}: YES (medium-rising-to-high under
cold-state-heavy workloads).

\begin{center}\rule{0.5\linewidth}{0.5pt}\end{center}

\subsubsection{\texorpdfstring{Q3.0-STM-013 ---
\texttt{drain\_exec}\textquotesingle s synthetic exec\_key uses
\texttt{origin\_lo,\ origin\_hi\ \&\ \textasciitilde{}kFlagMask} ---
adversarial origins force MVCC slot collisions
{[}Medium{]}}{Q3.0-STM-013 --- drain\_exec\textquotesingle s synthetic exec\_key uses origin\_lo, origin\_hi \& \textasciitilde kFlagMask --- adversarial origins force MVCC slot collisions {[}Medium{]}}}\label{q30-stm-013--drain_execs-synthetic-exec_key-uses-origin_lo-origin_hi--kflagmask--adversarial-origins-force-mvcc-slot-collisions-medium}

\textbf{Description}. \texttt{drain\_exec} derives the per-tx exec\_key
from:

\begin{verbatim}
ulong key_lo = ulong(v.origin_lo);
ulong key_hi = ulong(v.origin_hi & ~kFlagMask);
if (key_lo == 0UL && key_hi == 0UL) key_lo = 1UL;
\end{verbatim}

Two issues:

\begin{enumerate}
\def\labelenumi{\arabic{enumi}.}
\tightlist
\item
  Two different origins differing only in the top 2 bits
  (\texttt{kFlagMask\ =\ 0xC0000000}) produce \textbf{identical}
  exec\_keys. The flag bits are consumed; the rest is the key.
\item
  The "never empty" remap to \texttt{key\_lo=1} means every tx with
  origin=0 collides on the slot for key=(1, 0).
\end{enumerate}

\textbf{Affected}: \texttt{quasar\_wave.metal:902-905},
\texttt{quasar\_wave.cu:627-629}.

\textbf{Attack scenario}. Adversary submits 1024 transactions all with
\texttt{origin\ =\ 0} and \texttt{origin\_hi\ =\ 0xC0000000} (both flag
bits set). All have exec\_key = (1, 0). All hit the same MVCC slot.
Block-STM repair amplification = O(N²). Combined with Q3.0-STM-003 (no
repair bound), the round never finalizes.

\textbf{Fix}.

\begin{enumerate}
\def\labelenumi{\arabic{enumi}.}
\tightlist
\item
  Hash the origin (full 64 bits, no flag mask) plus tx\_index into a
  real 128-bit key.
\item
  Reserve a flag region in \texttt{last\_writer\_tx} instead of
  co-opting origin\_hi bits.
\end{enumerate}

\textbf{Must-fix pre-launch}: SHOULD be fixed --- but if not, must be
paired with Q3.0-STM-003\textquotesingle s repair bound to prevent DoS.

\begin{center}\rule{0.5\linewidth}{0.5pt}\end{center}

\subsubsection{\texorpdfstring{Q3.0-STM-014 --- Repair path discards
\texttt{er.incarnation} on requeue and resets \texttt{tx\_index\_seq} is
not reset between rounds in tests
{[}Medium{]}}{Q3.0-STM-014 --- Repair path discards er.incarnation on requeue and resets tx\_index\_seq is not reset between rounds in tests {[}Medium{]}}}\label{q30-stm-014--repair-path-discards-erincarnation-on-requeue-and-resets-tx_index_seq-is-not-reset-between-rounds-in-tests-medium}

\textbf{Description}. \texttt{drain\_repair}
(\texttt{quasar\_wave.metal:1013-1039}):

\begin{verbatim}
for (uint i = 0u; i < budget; ++i) {
    ExecResult er;
    if (!ring_try_pop(repair_hdr, repair_items, er)) break;
    VerifiedTx v;
    v.tx_index  = er.tx_index;
    v.admission = 0u;
    v.gas_limit = er.gas_used;
    ...
    if (!ring_try_push(exec_hdr, exec_items, v)) {
\end{verbatim}

The incarnation that \texttt{drain\_validate} bumped
(\texttt{er.incarnation\ +=\ 1u} at line 969) is passed to
\texttt{drain\_repair}, but drain\_repair packs only
\texttt{tx\_index,\ admission,\ gas\_limit,\ origin\_lo,\ origin\_hi}
into VerifiedTx --- \textbf{dropping the incarnation}. On the next exec,
drain\_exec writes \texttt{er.incarnation\ =\ 0u} (line 909). The
incarnation counter resets on every repair cycle.

This means LP-010 §"Repair monotonicity" ("each repair increments
incarnation; invalid incarnations cannot commit") cannot be enforced
because the kernel doesn\textquotesingle t carry incarnation through the
round-trip. A validator that wanted to enforce
\texttt{MAX\_TOTAL\_REPAIRS} (Q3.0-STM-003) couldn\textquotesingle t
even count repairs per tx.

\textbf{Affected}: \texttt{quasar\_wave.metal:1023-1029},
\texttt{quasar\_wave.cu:723-729} (same bug). \texttt{drain\_exec} line
909 (Metal) / 633 (CUDA) reset incarnation.

\textbf{Fix}. VerifiedTx layout needs an \texttt{incarnation} field, OR
Repair → Exec must travel through a different envelope (RepairTx with
full ExecResult including incarnation), OR incarnation must be stored
per tx\_index in a side table (FiberSlot already has the field but
it\textquotesingle s unused).

The FiberSlot has incarnation at offset 32 (\texttt{uint\ incarnation}
in the FiberSlot struct, line 261 in metal). drain\_exec writes 0 there
because FiberSlots are never indexed. Use the FiberSlot table:
drain\_validate increments \texttt{fibers{[}tx\_index{]}.incarnation};
drain\_exec reads it.

\textbf{Must-fix pre-launch}: YES (medium) --- without this,
Q3.0-STM-003 is unfixable.

\begin{center}\rule{0.5\linewidth}{0.5pt}\end{center}

\subsubsection{\texorpdfstring{Q3.0-STM-015 ---
\texttt{keccak256\_thread}/\texttt{keccak256\_local} final-block padding
works at boundaries; verified --- but uses LITTLE-ENDIAN lane absorption
which differs from Ethereum keccak (which is also little-endian, so this
is correct)
{[}Informational{]}}{Q3.0-STM-015 --- keccak256\_thread/keccak256\_local final-block padding works at boundaries; verified --- but uses LITTLE-ENDIAN lane absorption which differs from Ethereum keccak (which is also little-endian, so this is correct) {[}Informational{]}}}\label{q30-stm-015--keccak256_threadkeccak256_local-final-block-padding-works-at-boundaries-verified--but-uses-little-endian-lane-absorption-which-differs-from-ethereum-keccak-which-is-also-little-endian-so-this-is-correct-informational}

\textbf{Description}. I traced \texttt{keccak256} in both backends
carefully against FIPS 202 + Ethereum\textquotesingle s keccak (the
pre-NIST keccak with padding \texttt{0x01\ ...\ 0x80} in little-endian
byte ordering). Both backends absorb lanes via:

\begin{verbatim}
s[lane] ^= ulong(data[i]) << (i % 8u * 8u);
\end{verbatim}

which is little-endian byte-to-lane mapping. This is the Ethereum
convention. Cross-checked against \texttt{pyethereum/keccak.py}. This is
correct, BUT it must be cross-checked against any CPU reference
implementation used elsewhere in luxcpp (e.g.,
\texttt{lib/evm/state/processor.cpp}). If the CPU reference uses the
NIST SHA-3 byte-ordering (different padding: \texttt{0x06\ ...\ 0x80}),
the GPU keccak will diverge from CPU.

The \texttt{keccak256} implementations in Metal and CUDA agree
byte-for-byte. \textbf{No determinism risk between Metal and CUDA from
keccak.}

\textbf{Affected}: N/A --- verified correct.

\textbf{Recommendation}. Add a unit test that hashes a fixed input
through Metal keccak, CUDA keccak, and the CPU
\texttt{crypto::keccak256} used by the rest of cevm. Assert
byte-identical output.

\textbf{Must-fix pre-launch}: NO --- but the cross-impl test SHOULD land
for confidence.

\begin{center}\rule{0.5\linewidth}{0.5pt}\end{center}

\subsubsection{\texorpdfstring{Q3.0-STM-016 --- Ring requeue on
push-fail uses \texttt{ring\_try\_push(self\_hdr,\ ...)} which goes to
the TAIL; out-of-order ring rotation
{[}Low{]}}{Q3.0-STM-016 --- Ring requeue on push-fail uses ring\_try\_push(self\_hdr, ...) which goes to the TAIL; out-of-order ring rotation {[}Low{]}}}\label{q30-stm-016--ring-requeue-on-push-fail-uses-ring_try_pushself_hdr--which-goes-to-the-tail-out-of-order-ring-rotation-low}

\textbf{Description}. When \texttt{drain\_validate} fails to push to
commit (line 1002):

\begin{verbatim}
(void)ring_try_push(validate_hdr, validate_items, er);
\end{verbatim}

This pushes the failed tx to the \emph{tail} of the validate ring.
Subsequent pops will see other items first, then this one. This means
the \texttt{er} whose MVCC was already mutated (Q3.0-STM-001) lands
behind newer execrequests in the same tick. The order ExecResult items
were popped is NOT preserved.

Combined with Q3.0-STM-001, this means the order of validate pop on the
next tick is timing-dependent.

\textbf{Affected}: lines 938-940 (drain\_exec), 970-972 (drain\_validate
to repair), 1001-1002 (drain\_validate to commit), 1031-1033
(drain\_repair), 1056-1059 (drain\_state\_resp). Same pattern in CUDA
backend.

\textbf{Fix}. Don\textquotesingle t requeue to tail. Either:

\begin{enumerate}
\def\labelenumi{\arabic{enumi}.}
\tightlist
\item
  Stop popping when downstream is full (already partially done ---
  \texttt{break} on push fail). The pop\textquotesingle d item is then
  lost. Need to NOT pop until downstream has space.
\item
  Use a peek-first-then-pop pattern: query the downstream
  ring\textquotesingle s free space, only pop if there\textquotesingle s
  space.
\end{enumerate}

The cleanest approach: peek tail-vs-head before popping. Standard SPSC
ring pattern.

\textbf{Must-fix pre-launch}: NO (low). Doesn\textquotesingle t directly
cause divergence (the tx is still in the ring), but interacts badly with
Q3.0-STM-001 and adds confusion.

\begin{center}\rule{0.5\linewidth}{0.5pt}\end{center}

\subsection{Cross-Cutting Architectural
Concerns}\label{cross-cutting-architectural-concerns}

\subsubsection{LP-vs-Implementation Gap}\label{lp-vs-implementation-gap}

LP-010 specifies a 3-tier validation pipeline (lane clocks, MVCC,
semantic reducers), commit horizons, lane refraction, deterministic
contention manager, hard livelock limits, and multi-GPU sharding.
\textbf{None of these are implemented.} The kernel ships flat MVCC +
repair-only + finalization-on-empty. Calling this Block-STM 3.0 in
production by 2025-12-25 is misleading.

The LP itself blurs this --- §"Implementation Plan" lists 3.0 features
at v0.42-v0.49, with v0.49 being "Formal CPU reference (4.0)". Per the
LP, this is post-launch work. \textbf{But the LP\textquotesingle s
"Correctness Invariants" §1-7 are stated as MUST hold for 3.0.} They
cannot hold without the implementation.

\textbf{Recommendation}: Either rename the launch artifact to "QuasarSTM
2.5 / GPU Block-STM" and ship it honestly, OR delay launch until v0.45+.
Don\textquotesingle t ship v0.40 calling itself 3.0.

\subsubsection{Forbidden Patterns Audit}\label{forbidden-patterns-audit}

LP-010 §"Forbidden Patterns" says QuasarSTM must not use:

\begin{itemize}
\tightlist
\item
  one global STM clock --- N/A (not used)
\item
  one global version map lock --- N/A
\item
  retry-until-success GPU loops --- VIOLATED at \texttt{ring\_try\_pop}
  (Metal line 298-311, CUDA line 268-281). The CAS loop is unbounded in
  principle. In practice it terminates because there are finitely many
  concurrent threads, but the pattern is exactly what the LP forbids.
\item
  nondeterministic conflict victim choice --- VIOLATED via Q3.0-STM-001
  (timing decides who repairs)
\item
  opcode-level STM --- N/A (it\textquotesingle s tx-level, correct)
\item
  CPU conflict manager --- partially violated via host\textquotesingle s
  \texttt{max\_epochs} choice (Q3.0-STM-003)
\item
  host-side repair scheduler --- N/A
\item
  per-version \texttt{malloc} --- N/A (uses arena)
\item
  global hot-key spinlocks --- N/A
\end{itemize}

\textbf{Two of nine forbidden patterns are present} in the v0.40 kernel.

\begin{center}\rule{0.5\linewidth}{0.5pt}\end{center}

\subsection{Blue Handoff}\label{blue-handoff}

\subsubsection{What Blue got right}\label{what-blue-got-right}

\begin{itemize}
\tightlist
\item
  The LP architecture (3-tier validation, lane refraction, commit
  horizons, semantic reducers) is solid and reflects current research.
  The mistake is shipping before implementing it.
\item
  Wave-tick scheduler with bounded budgets --- correct anti-livelock
  approach at the dispatch layer (kernel exits, GPU yields).
\item
  \texttt{RingHeader} layout and the relaxed-atomic +
  \texttt{mem\_device} fence pattern is correct \emph{if applied with
  proper memory ordering} (it isn\textquotesingle t --- see
  Q3.0-STM-005).
\item
  Inline keccak verified correct, byte-identical between Metal and CUDA.
\item
  \texttt{static\_assert} on layout sizes catches host/device drift.
\item
  Basic determinism test (\texttt{test\_root\_determinism}) exists and
  passes for the trivial 16-tx case.
\end{itemize}

\subsubsection{What Blue missed}\label{what-blue-missed}

\begin{itemize}
\tightlist
\item
  No cross-backend determinism test harness.
\item
  No CPU reference for differential fuzzing (LP-010 v0.49 promises this;
  not landed).
\item
  No adversarial workload tests (1024 same-key txs, hash-collision keys,
  multi-cold-page txs, push-during-close).
\item
  No real signature verification on the consensus path --- both backends
  ship placeholder verifiers that are trivially forgeable.
\item
  No state\_root computation --- block\_hash is uncorrelated with state.
\item
  No anti-livelock enforcement despite the LP marking it "mandatory on
  GPU."
\end{itemize}

\subsubsection{Fix priority for Blue
(ordered)}\label{fix-priority-for-blue-ordered}

\begin{enumerate}
\def\labelenumi{\arabic{enumi}.}
\tightlist
\item
  \textbf{Q3.0-STM-008} --- CUDA verify\_signature\_stub. One-line
  forgery. Replace with the same HMAC-keccak path as Metal, OR
  (preferred) gate the round-finalization on host-side signature
  verification using \texttt{luxfi/crypto}.
\item
  \textbf{Q3.0-STM-009} --- hardcoded master secret. Rip out the
  HMAC-keccak verifier entirely; either ship real BLS/Ringtail/MLDSA by
  2025-12-25 or defer GPU vote aggregation and run host-side signature
  verification.
\item
  \textbf{Q3.0-STM-001} --- commit-ring backpressure. Refactor
  \texttt{drain\_validate} so the MVCC mutation is contingent on
  commit-push success.
\item
  \textbf{Q3.0-STM-002} --- non-atomic key claim in mvcc\_locate. Use
  atomic CAS on a sentinel field with proper memory ordering.
\item
  \textbf{Q3.0-STM-004} --- port v0.40 DAG to CUDA; build cross-backend
  test harness; add CPU reference.
\item
  \textbf{Q3.0-STM-003} + \textbf{Q3.0-STM-014} --- implement repair
  monotonicity (incarnation in FiberSlot) and enforce
  MAX\_TOTAL\_REPAIRS with deterministic abort.
\item
  \textbf{Q3.0-STM-005} --- fix memory-order on DAG state atomics
  (release/acquire instead of relaxed).
\item
  \textbf{Q3.0-STM-006} --- move tx\_index\_seq increment after
  decode-push success.
\item
  \textbf{Q3.0-STM-007} --- eliminate the vote\_verified bitmap; verify
  inline in drain\_vote.
\item
  \textbf{Q3.0-STM-011} --- implement state\_root accumulation.
\item
  \textbf{Q3.0-STM-010} --- gate push\_txs on closing\_flag.
\item
  \textbf{Q3.0-STM-012} --- implement loaded sentinel; coalesce
  multi-page suspends.
\item
  \textbf{Q3.0-STM-013} --- replace synthetic exec\_key derivation with
  hash of full origin.
\item
  \textbf{Q3.0-STM-016} --- peek-then-pop ring pattern.
\item
  \textbf{Q3.0-STM-015} --- add cross-impl keccak test (informational).
\end{enumerate}

\subsubsection{Re-review scope}\label{re-review-scope}

Re-test required after fixes for:

\begin{itemize}
\tightlist
\item
  Cross-backend determinism harness (Metal vs CUDA vs CPU reference) ---
  full 19-binary GPU test surface plus new adversarial workloads (1024
  same-key, FNV-collision keys, multi-cold-page, push-during-close,
  race-during-DAG).
\item
  Repair-bounded-abort behavior (1024 same-key + bounded
  MAX\_TOTAL\_REPAIRS = deterministic block rejection).
\item
  State\_root accumulation produces stable digest under reordered
  execution within Block-STM constraints.
\item
  Real signature verifiers in all three lanes --- must use
  \texttt{luxfi/crypto} and reject the test suite\textquotesingle s
  HMAC-keccak forgeries.
\end{itemize}

\begin{center}\rule{0.5\linewidth}{0.5pt}\end{center}

\subsection{Recommendation}\label{recommendation}

\textbf{do-not-ship} as Block-STM 3.0 by 2025-12-25.

The v0.40 substrate can ship as \textbf{QuasarSTM 2.5 / GPU Block-STM}
(honest naming) once the must-fix-pre-launch items above are addressed.
The LP-010 3.0 features (lanes, three-tier validation, commit horizons,
deterministic contention manager) are real engineering work --- months,
not weeks --- and shipping a kernel that lacks them while claiming the
3.0 invariants is a credibility risk and an attack surface.

Critical \& high findings (Q3.0-STM-001, 002, 003, 004, 005, 006, 007,
008, 009) \textbf{MUST land before any mainnet round.} They are all
bounded local fixes; none requires the 3.0 rewrite. Estimated total
work: 2-3 weeks if focused, plus one week for cross-backend test harness
and one week for differential fuzzing against a CPU reference.

If those land and pass cross-backend determinism on the adversarial
workload suite, this can ship --- but as 2.5, not 3.0.

\begin{center}\rule{0.5\linewidth}{0.5pt}\end{center}

RED COMPLETE. Findings recorded. Total: 4 critical, 5 high, 4 medium, 3
low/informational

Top 3 for Blue to fix:

\begin{enumerate}
\def\labelenumi{\arabic{enumi}.}
\tightlist
\item
  Q3.0-STM-008 --- CUDA byte-equality "verifier" forges any signature
  trivially
\item
  Q3.0-STM-001 --- commit-ring backpressure produces nondeterministic
  repair counts → root divergence
\item
  Q3.0-STM-002 --- non-atomic MVCC key claim allows torn reads +
  double-claim → root divergence
\end{enumerate}

Re-review needed: yes --- cross-backend determinism harness +
adversarial workload suite + real signature verifiers Recommendation:
do-not-ship as 3.0; ship as 2.5 after the 9 critical+high items land

\end{document}
